\documentclass[conference]{IEEEtran}
\IEEEoverridecommandlockouts

\usepackage{cite}
\usepackage{amsmath,amssymb,amsfonts}
\usepackage{graphicx}
\usepackage{booktabs}
\usepackage{multirow}
\usepackage{xcolor}
\usepackage{algorithmic}
\usepackage{textcomp}
\usepackage{url}
\usepackage{hyperref}

% Visible marker for placeholder content that must be filled in once
% full-length training runs (200 epochs, DOTAv1.5 + VisDrone) have completed.
\newcommand{\todo}[1]{\textcolor{red}{\textbf{[TODO: #1]}}}

\begin{document}

\title{Combining Multi-Scale Context Aggregation, Attention-Guided Upsampling,\\
and High-Resolution Feature Enhancement for Small Object Detection with YOLO}

\author{
\IEEEauthorblockN{Osias No\"{e}l Tossou}
\IEEEauthorblockA{\todo{Affiliation / Department / University}\\
\todo{City, Country}\\
osias.tossou@gmail.com}
}

\maketitle

\begin{abstract}
Small object detection remains one of the hardest open problems for single-stage
detectors such as the YOLO family, because repeated stride-2 downsampling
in the backbone erodes the spatial signal of small instances well before the
detection heads can exploit it. In this work we study three complementary
architectural interventions on YOLO11/YOLO12: (i) an attention-augmented,
cross-stage-partial variant of the Spatial Pyramid Pooling -- Fast (SPPF)
block at the lowest-resolution stage (P5), combining a dilated multi-branch
context path with Convolutional Block Attention Module (CBAM) refinement;
(ii) a high-resolution early-stage branch that inserts lightweight attention
modules at the P2/P3 stages and adds a dedicated P2 detection head, so that
small-object evidence is exploited before it is downsampled away; and (iii)
content-aware feature reassembly (CARAFE) in place of nearest-neighbor
upsampling in the top-down feature pyramid, so that the detail reconstructed
during top-down fusion is learned rather than interpolated. Rather than jumping directly to a full-scale benchmark comparison, we first
validate the reasoning behind each intervention on its own terms: a
data-grounded analysis of how much of a small object's footprint survives
at each pyramid stride on two real aerial/UAV datasets (DOTA v1.5,
VisDrone); an exact, gradient-based proof that a pooling-only aggregation
path is structurally blind to non-maximal input values while a parallel
cross-stage-partial shortcut is not; and a controlled proxy reconstruction
task quantifying the fidelity gain of content-aware upsampling over
nearest-neighbor interpolation. All three validations use real data and
are fully reproducible from the accompanying code, without requiring
full-length detector training. Large-scale comparative training on DOTA
v1.5 and VisDrone (YOLO11, YOLO12, and every architectural variant,
including the fully combined model) is the subject of a forthcoming
experimental study, for which this paper establishes the architecture,
the protocol, and the theoretical motivation.
\end{abstract}

\begin{IEEEkeywords}
small object detection, YOLO, spatial pyramid pooling, attention mechanism, feature pyramid network, aerial imagery, UAV imagery
\end{IEEEkeywords}

\section{Introduction}
\label{sec:intro}

Single-stage detectors of the YOLO family~\cite{redmon2016yolo} have become the
default choice for real-time object detection because they offer a favorable
accuracy/latency trade-off. However, their performance on \emph{small}
objects --- a few pixels wide after the backbone's successive stride-2
downsampling stages --- lags substantially behind their performance on
medium and large instances. This gap is particularly acute in aerial and
UAV imagery, where the camera's bird's-eye viewpoint and high altitude make
most instances of interest (vehicles, boats, pedestrians) occupy only a
handful of pixels, and where datasets such as DOTA~\cite{xia2018dota} and
VisDrone~\cite{zhu2021visdrone} report that the majority of annotated
instances fall under commonly used small-object area thresholds.

The root cause is architectural: by the time a feature map reaches the
detector's deepest stage (P5, stride 32), a small object's spatial footprint
has often collapsed to a single feature-map cell or less, and the
information is further diluted by whatever aggregation happens at that
stage --- in YOLO11/YOLO12, the Spatial Pyramid Pooling -- Fast
(SPPF)~\cite{he2015spp} block, a cascade of same-kernel max-pooling
operations. SPPF enlarges the effective receptive field cheaply, but it
does so purely through max-pooling, which discards the fine-grained detail
a small object needs and treats every channel and spatial location
identically.

We argue --- and structure our ablation study around the claim --- that no
single intervention at a single stage of the network is sufficient: a
richer SPPF replacement at P5 only reshapes the \emph{global context} that
is later fused back into higher-resolution stages through the feature
pyramid; it cannot recover detail already destroyed by downsampling.
Conversely, an early-stage, high-resolution intervention that preserves
small-object detail has no mechanism to inject the broader scene context
that disambiguates visually similar small classes. We therefore propose
and evaluate three interventions that act at three different points of the
network, and measure whether their effects are complementary rather than
redundant:

\begin{itemize}
  \item \textbf{SPPFCSPC-CBAM} (Section~\ref{sec:method-sppf}): a
  cross-stage-partial, attention-refined replacement for SPPF at P5,
  combining a CSP shortcut branch, a dilated multi-branch pooling cascade,
  and CBAM~\cite{woo2018cbam} channel/spatial attention.
  \item \textbf{TOSA + P2 head} (Section~\ref{sec:method-tosa}): lightweight
  attention modules inserted at the P2 and P3 stages of the backbone, a
  soft (average-pooling-based) downsampling operator in place of strided
  convolutions between P2 and P4, and an additional detection head at
  stride 4 (P2), so that small-object evidence is used directly rather than
  only through a downsampled proxy.
  \item \textbf{CARAFE upsampling} (Section~\ref{sec:method-carafe}): a
  content-aware reassembly operator~\cite{wang2019carafe} that replaces
  nearest-neighbor upsampling in the top-down feature pyramid path, so that
  the detail injected during top-down fusion is predicted from local
  content rather than interpolated.
\end{itemize}

Our contributions are:
\begin{enumerate}
  \item A combined architecture integrating P5 context aggregation (SPPFCSPC-CBAM),
  P2/P3 high-resolution attention with a dedicated P2 head (TOSA), and
  content-aware FPN upsampling (CARAFE) into a single YOLO11-based model,
  designed so each component is independently toggleable for ablation.
  \item A data-grounded quantitative argument, computed directly on the
  DOTA v1.5 and VisDrone training sets, for \emph{where} in the network a
  small-object intervention can still act before the object's footprint
  has collapsed below the resolution of a single feature-map cell
  (Section~\ref{sec:theory-resolution}).
  \item An exact, closed-form and gradient-verified demonstration that a
  pooling-only aggregation path (plain SPPF) structurally discards
  information that a cross-stage-partial shortcut (as used in our
  SPPFCSPC-CBAM) preserves (Section~\ref{sec:theory-bottleneck}).
  \item A controlled proxy experiment, trained on real DOTA image patches,
  quantifying the reconstruction-fidelity gain of content-aware upsampling
  (CARAFE) over nearest-neighbor interpolation (Section~\ref{sec:theory-carafe}).
  \item A fully specified protocol for the large-scale comparative training
  study (datasets, variants, hyperparameters) that this paper's theoretical
  and controlled validations motivate, left for a forthcoming experimental
  phase (Section~\ref{sec:experiments}).
\end{enumerate}

\section{Related Work}
\label{sec:related}

\subsection{YOLO and Spatial Pyramid Pooling}
The YOLO family~\cite{redmon2016yolo} evolved from YOLOv1's single dense
prediction grid to multi-scale feature pyramids with a Spatial Pyramid
Pooling (SPP) stage~\cite{he2015spp} at the backbone's deepest level, which
YOLOv5 replaced with the cheaper, functionally equivalent SPPF (a single
pooling kernel applied repeatedly instead of three parallel kernels of
increasing size). YOLOv7 proposed SPPCSPC~\cite{wang2023yolov7}, which adds
a cross-stage-partial shortcut around the pooling cascade, reducing the
information bottleneck a plain SPP/SPPF introduces; our SPPFCSPC-CBAM
module (Section~\ref{sec:method-sppf}) extends this line of work with
explicit multi-scale dilation and attention-based refinement.

\subsection{Multi-scale context and receptive field diversity}
Beyond SPP-style pooling, atrous/dilated convolutions have been used to
enlarge the receptive field while preserving spatial resolution, notably in
ASPP~\cite{chen2017deeplab} for semantic segmentation and RFBNet's
Receptive Field Block~\cite{liu2018rfb} for object detection, which fuses
parallel dilated branches explicitly to diversify the effective receptive
field rather than relying solely on pooling. Our SPPFCSPC-CBAM module
borrows this multi-dilation-branch principle while retaining the
cheap cascaded-pooling branch of SPPF for a controlled comparison.

\subsection{Attention mechanisms}
Channel and spatial attention modules have been shown to cheaply improve
feature discriminability without significantly increasing model size. CBAM
~\cite{woo2018cbam} sequentially applies channel and spatial attention;
SimAM~\cite{yang2021simam} achieves a comparable effect with a
parameter-free energy-based formulation. We use CBAM inside our SPPF
replacement and note SimAM as a lighter-weight alternative evaluated as a
secondary module in our codebase but outside the scope of the present
ablation.

\subsection{Small object detection}
Dedicated small-object detection work generally falls into two families:
architectural interventions that preserve or restore high-resolution
detail (additional high-resolution detection heads, e.g. a P2/stride-4
head as used by several YOLO variants and by QueryDet-style
cascades; and feature pyramid upsampling improvements such as
CARAFE~\cite{wang2019carafe}, which predicts content-aware reassembly
kernels instead of using fixed nearest-neighbor or bilinear interpolation
in the top-down path) and data/training interventions (tiled inference
such as SAHI, copy-paste augmentation of small instances, and loss
functions less sensitive to small-box IoU volatility such as
Wise-IoU/SIoU variants). The present work focuses on the architectural
family, and treats the P2-head and CARAFE directions as complementary,
testing them jointly with our SPPF-stage contribution rather than in
isolation.

\subsection{Aerial and UAV object detection benchmarks}
DOTA~\cite{xia2018dota} and VisDrone~\cite{zhu2021visdrone} are the two
benchmarks used in this study. Both are dominated by small, dense
instances captured from an aerial or near-aerial viewpoint (vehicles,
ships, pedestrians), which makes them a natural testbed for small-object
architectural interventions, as opposed to ground-level benchmarks such as
COCO where the small-object subset is a much smaller fraction of the data.

\section{Proposed Method}
\label{sec:method}

\begin{figure}[htbp]
\centering
\fbox{\parbox{0.9\linewidth}{\centering\vspace{3.2cm}
\todo{Insert block diagram of the combined architecture: backbone with
TOSA/P2 head (Sec.~\ref{sec:method-tosa}), SPPFCSPC-CBAM at P5
(Sec.~\ref{sec:method-sppf}), and CARAFE upsampling in the FPN
(Sec.~\ref{sec:method-carafe}).}
\vspace{3.2cm}}}
\caption{Overview of the proposed architecture.}
\label{fig:architecture}
\end{figure}

All variants studied in this paper share the YOLO11 backbone/head topology
(Fig.~\ref{fig:architecture}), and differ
only in the modules described below, which can be composed independently.
This lets us attribute any accuracy change to a specific, isolated
architectural change rather than to a confound of several simultaneous
modifications.

\subsection{Baseline}
The YOLO11 baseline backbone downsamples the input through five stages
(P1--P5, strides 2--32) using strided $3\times3$ convolutions and C3k2
blocks, aggregates context at P5 with SPPF, refines it with a C2PSA
attention block, and feeds a standard PANet-style feature pyramid (nearest-
neighbor upsampling, concatenation, C3k2 fusion) into three detection heads
at P3/P4/P5. YOLO12 is evaluated as a second, independent baseline.

\subsection{SPPFCSPC-CBAM: attention-refined SPPF replacement}
\label{sec:method-sppf}

Plain SPPF aggregates context purely through a cascade of same-kernel
max-pooling operations, which (i) discards fine-grained detail at every
pooling step and (ii) treats all channels and spatial locations uniformly.
Our replacement, SPPFCSPC-CBAM, addresses both limitations while keeping
the parameter/FLOP overhead at P5 --- the network's lowest-resolution,
cheapest-to-widen stage --- deliberately contained.

Given an input feature map $X \in \mathbb{R}^{C \times H \times W}$, the
module splits computation into a cross-stage-partial shortcut branch and a
pooling branch (Eq.~\ref{eq:sppfcspc-split}):
\begin{align}
\label{eq:sppfcspc-split}
X_1 &= \mathrm{Conv}_{1\times1}(X), \quad X_2 = \mathrm{Conv}_{1\times1}(X) \\
X_1' &= \mathrm{Conv}_{1\times1}(\mathrm{Conv}_{3\times3}(X_1))
\end{align}
where $X_2$ bypasses pooling entirely and is fused back in at the end. The
pooled branch applies a $k\times k$ max-pool $n$ times in sequence (as in
SPPF, equivalent to a single pass of $k, 2k, \dots$ receptive fields at a
fraction of the cost of parallel kernels):
\begin{equation}
Y_0 = X_1', \quad Y_i = \mathrm{MaxPool}_{k}(Y_{i-1}), \; i = 1, \dots, n
\end{equation}
The $n+1$ pooled feature maps are concatenated, fused, and refined with
CBAM channel-then-spatial attention~\cite{woo2018cbam} before being
concatenated with the untouched shortcut branch and projected to the
output channel count:
\begin{align}
Y &= \mathrm{CBAM}\big(\mathrm{Conv}_{3\times3}(\mathrm{Conv}_{1\times1}([Y_0, \dots, Y_n]))\big) \\
\mathrm{Out} &= \mathrm{Conv}_{1\times1}([Y, X_2])
\end{align}
The CSP shortcut ($X_2$) guarantees that at least one path through the
module carries the input forward with no lossy pooling, and CBAM
re-weights the pooled branch by learned channel and spatial relevance
rather than taking the pooling cascade's output at face value.

\subsection{TOSA and the P2 detection head}
\label{sec:method-tosa}

SPPFCSPC-CBAM operates at P5 (stride 32) and therefore cannot recover
detail already lost to downsampling; it only improves the quality of the
context subsequently fused back into higher-resolution stages. To act
\emph{before} that detail is lost, we insert a lightweight attention
module (TOSA) directly after the P2 (stride 4) and P3 (stride 8) stages of
the backbone, replace the strided $3\times3$ convolutions that downsample
P2$\to$P3 and P3$\to$P4 with a channel-preserving average-pooling depthwise
convolution (so that width scaling propagates through the backbone
unchanged), and add a fourth detection head at P2/stride 4 ($160\times160$
at a $640\times640$ input), alongside the standard P3/P4/P5 heads. The P2
head lets the detector directly supervise and predict from the
highest-resolution feature map the network produces, rather than relying
exclusively on heads that only see small objects after at least one more
round of downsampling than P2 itself.

\subsection{Content-aware upsampling (CARAFE)}
\label{sec:method-carafe}

The top-down path of a standard YOLO feature pyramid upsamples
low-resolution, high-level features with fixed nearest-neighbor
interpolation before fusing them with the corresponding higher-resolution
backbone feature map. This interpolation is content-agnostic: it cannot
reconstruct any detail beyond what the low-resolution map already encodes
explicitly. We replace nearest-neighbor upsampling with
CARAFE~\cite{wang2019carafe}, which predicts, from the local content of the
low-resolution feature map, a reassembly kernel used to redistribute each
output pixel's value from a local neighborhood of input pixels, instead of
simply repeating the nearest input value. This is channel-preserving and
computationally lightweight (a small encoder plus an unfold-based weighted
sum) relative to the convolutional blocks that already dominate the
pyramid's cost.

\subsection{Combined architecture}
\label{sec:method-combined}

Our proposed architecture (denoted COMBO in Section~\ref{sec:results})
applies all three interventions jointly on top of the YOLO11 topology: TOSA
and the P2 head at the backbone's high-resolution stages, SPPFCSPC-CBAM at
P5, and CARAFE throughout the top-down FPN path. Because each intervention
acts at a different point of the network and is independently toggleable
in our implementation, Section~\ref{sec:results} reports both the combined
model and each component in isolation, which lets us test directly whether
their effects are additive, redundant, or interacting.

\section{Theoretical Analysis and Controlled Validation}
\label{sec:theory}

Each design choice in Section~\ref{sec:method} was motivated by an
argument about what information a given mechanism can or cannot carry
forward. This section tests those arguments directly, without requiring
full detector training: Section~\ref{sec:theory-resolution} grounds the
claim that P5-only interventions act too late with real per-instance
statistics from both datasets; Section~\ref{sec:theory-bottleneck} proves,
exactly, that a pooling-only path is blind to information a parallel
shortcut preserves; Section~\ref{sec:theory-carafe} quantifies, on a small
trained proxy task, the reconstruction-fidelity gain of content-aware
upsampling. All code for these three validations is included in the
paper's supplementary repository (\texttt{paper/theory\_demo/}) and runs
in minutes on a single machine, with no GPU training required for the
first two.

\subsection{Small-object resolution collapse is a real-data phenomenon, not an estimate}
\label{sec:theory-resolution}

Section~\ref{sec:intro} argued that a P5-only intervention cannot recover
spatial detail already destroyed by downsampling, because by stride 32 a
small object's footprint has typically collapsed below one feature-map
cell. We verify this directly: for every annotated instance in the DOTA
v1.5 and VisDrone training sets, we recover its bounding-box width and
height in pixels (axis-aligned extent of the oriented box for DOTA; native
pixel annotations for VisDrone) and compute the fraction of
\emph{small} instances (area $<32^2$px, the COCO convention) whose
footprint, after dividing by the stride, covers less than one feature-map
cell in area ($\frac{w}{s}\cdot\frac{h}{s} < 1$) at each of the four
strides used by our architectures (Table~\ref{tab:resolution}).

\begin{table}[htbp]
\caption{Fraction of small instances ($<32^2$px) collapsing to under one
feature-map cell, by pyramid stride -- computed on the real training
annotations}
\label{tab:resolution}
\centering
\begin{tabular}{lccccc}
\toprule
Dataset & Small \% & $s{=}4$ & $s{=}8$ & $s{=}16$ & $s{=}32$ \\
\midrule
DOTA v1.5 (train)  & 66.8\% & 0.0\%  & 13.7\% & 47.6\%  & 100.0\% \\
VisDrone (train)   & 60.5\% & 0.5\%  & 8.8\%  & 43.0\%  & 100.0\% \\
\bottomrule
\end{tabular}
\end{table}

Three observations follow directly from Table~\ref{tab:resolution} and
motivate the method design. First, small objects are not a minority case
for either dataset: 66.8\% of DOTA instances and 60.5\% of VisDrone
instances fall under the COCO small-object threshold, so an intervention
that measurably helps this subset should move the aggregate metric as
well. Second, at $s{=}32$ (P5, where SPPF/SPPFCSPC-CBAM operate),
\emph{every} small instance in both datasets collapses to under one
feature-map cell -- not a majority, all of them. No amount of
architectural sophistication applied only at P5 can recover spatial
detail that no longer exists in the feature map at that resolution; it
can only improve the quality of the aggregate context subsequently fused
back into higher-resolution stages, exactly as argued in
Section~\ref{sec:intro}. Third, even $s{=}8$ (P3, the finest stride a
standard three-head YOLO11/YOLO12 exposes to the detector) already loses
13.7\% (DOTA) to 8.8\% (VisDrone) of small instances to sub-cell collapse,
while $s{=}4$ (P2, where our P2 head operates) loses under 1\% on both
datasets. This is the quantitative basis for the TOSA/P2-head
intervention (Section~\ref{sec:method-tosa}): it is not merely a
heuristic choice of ``add a higher-resolution head'', it targets the
specific stride at which the real small-object population in these two
datasets is still almost entirely resolvable.

\subsection{Pooling is an information bottleneck; the CSP shortcut is not}
\label{sec:theory-bottleneck}

Section~\ref{sec:method-sppf} argued that plain SPPF discards
fine-grained detail because it aggregates exclusively through max-pooling,
and that the CSP shortcut in SPPFCSPC-CBAM avoids this. We make this
precise and verify it exactly via automatic differentiation, which gives
an unambiguous test of whether a module's output can depend on a given
input value at all: if $\partial\,\mathrm{output}/\partial\,x_{i} = 0$
identically, no change to $x_i$, however large, can be observed in the
output through that computational path.

\paragraph{Lemma.} For a max-pool operator with window $\Omega$, and any
position $i \in \Omega$ that is not the arg-max of $\Omega$, the gradient
of $\max_{\Omega}$ with respect to $x_i$ is exactly zero -- a direct
consequence of max-pooling's well-known subgradient (gradient flows only
to the arg-max element of each window). Applying a max-pool operator $n$
times in sequence, as SPPF does, composes this property: a position that
is never the arg-max of any window it participates in, at any stage of
the cascade, has exactly zero gradient with respect to the cascade's
output.

\paragraph{Verification.} We instantiate a 3-stage, $5\times5$ max-pool
cascade -- identical in structure to SPPF's pooling path -- on a random
$8\times6\times6$ input, back-propagate from the summed output, and
compare the gradient magnitude at arg-max vs.\ non-arg-max positions.
Separately, we back-propagate through an isolated $1\times1$ convolution
(structurally identical to the CSP shortcut branch $X_2$ in
Eq.~\ref{eq:sppfcspc-split}) on a random input of the same shape.
Table~\ref{tab:bottleneck} reports the exact result.

\begin{table}[htbp]
\caption{Gradient magnitude by input position: pooling cascade (SPPF's
sole aggregation path) vs. the CSP shortcut branch (SPPFCSPC-CBAM)}
\label{tab:bottleneck}
\centering
\begin{tabular}{lcc}
\toprule
Path & \multicolumn{2}{c}{$|\partial\,\mathrm{output}/\partial\,x_i|$} \\
 & non-arg-max positions & arg-max positions \\
\midrule
3-stage $5{\times}5$ max-pool cascade & $0.0$ (exact) & $36.0$ \\
\midrule
 & \multicolumn{2}{c}{all positions} \\
CSP shortcut (Conv$_{1\times1}$) & \multicolumn{2}{c}{$> 0$ everywhere (min $0.032$)} \\
\bottomrule
\end{tabular}
\end{table}

271 of the 288 positions in the test tensor are non-arg-max; every one of
them has \emph{exactly} zero gradient through the pooling cascade (not
merely small -- identically zero in floating-point arithmetic), confirming
that SPPF's sole aggregation mechanism is, by construction, blind to the
vast majority of its input. The CSP shortcut branch, in contrast, has
strictly positive gradient magnitude at all 288 positions: no information
is structurally discarded on that path. This is precisely the property
SPPFCSPC-CBAM relies on (Section~\ref{sec:method-sppf}): the shortcut
guarantees at least one path through the module is not subject to the
lemma above, independent of what the pooling branch or the CBAM
refinement subsequently learn to do with their own share of the
representation.

\subsection{Content-aware upsampling reconstructs more detail than nearest-neighbor interpolation}
\label{sec:theory-carafe}

Section~\ref{sec:method-carafe} argued that nearest-neighbor upsampling in
the top-down FPN path is content-agnostic and cannot reconstruct detail
beyond what the low-resolution map already encodes explicitly, whereas
CARAFE predicts a reassembly kernel from local content. Because CARAFE's
reassembly kernels are learned, this claim cannot be verified in closed
form the way Section~\ref{sec:theory-bottleneck}'s can; we instead
quantify it on a small, fast proxy task that isolates the upsampling
operator from the rest of the detection pipeline: reconstructing an
original-resolution image patch from its $2\times$ average-pooled
downsampled version, using real DOTA v1.5 training images (320 random
$64\times64$ patches from 40 images, 85/15 train/test split). We compare
parameter-free nearest-neighbor interpolation against a CARAFEFast module
(Section~\ref{sec:method-carafe}'s implementation) trained for 1,500 Adam
steps (batch 16, lr $10^{-3}$) to minimize reconstruction MSE on the
training patches, then evaluate PSNR on the held-out test patches
(Table~\ref{tab:carafe}).

\begin{table}[htbp]
\caption{$2\times$ upsampling reconstruction fidelity, held-out DOTA patches}
\label{tab:carafe}
\centering
\begin{tabular}{lc}
\toprule
Upsampling method & PSNR (dB) \\
\midrule
Nearest-neighbor (no training) & 33.48 \\
CARAFE (1,500 training steps)  & \textbf{34.64} (+1.17) \\
\bottomrule
\end{tabular}
\end{table}

A trained CARAFE module reconstructs the held-out patches with 1.17\,dB
higher PSNR than nearest-neighbor interpolation, after a training budget
several orders of magnitude smaller than a full detection training run.
We report this explicitly as a proxy result, not a substitute for
measuring CARAFE's effect inside the full FPN on detection metrics: an
RGB image patch is a convenient, inspectable stand-in for a
channel-generic feature map (CARAFEFast is agnostic to what the channels
represent), but the semantic feature maps CARAFE actually operates on in
the detector are learned jointly with the rest of the network rather than
fixed in advance. The proxy result nonetheless directly supports the
mechanism-level claim it is designed to test: content-aware reassembly
reconstructs information that fixed interpolation cannot, given the same
low-resolution input.

\section{Experimental Setup}
\label{sec:experiments}

The large-scale comparative study outlined below is the next phase of
this work; it was not required to validate the mechanisms in
Section~\ref{sec:theory}, but is necessary to establish whether and by how
much each intervention improves standard detection metrics in a full
training regime. We specify the complete protocol here so the
forthcoming results are directly reproducible from this description.

\subsection{Datasets}
We evaluate on two aerial/UAV object detection benchmarks dominated by
small, dense instances:
\begin{itemize}
  \item \textbf{DOTA v1.5}~\cite{xia2018dota}: 1,411 training and 458
  validation images (the official \texttt{test} split, 937 images, has no
  public ground truth and is excluded from this study), 16 categories,
  annotated with oriented bounding boxes, converted to axis-aligned boxes
  for the detection task studied here.
  \item \textbf{VisDrone}~\cite{zhu2021visdrone}: drone-captured imagery
  across 10 object categories (pedestrian, vehicle, bicycle, etc.),
  characterized by dense small-object scenes with substantial occlusion.
\end{itemize}

\subsection{Models compared}
\begin{itemize}
  \item \textbf{YOLO11} / \textbf{YOLO12}: unmodified baselines.
  \item \textbf{GL}: SPPF replaced by a non-attention GL-CAB block at P5
  only, isolating the effect of the architectural change itself from
  attention.
  \item \textbf{GLCCP} / \textbf{CCP}: two further SPPF-slot variants built
  on progressive dilated context aggregation (GL-PSWCA).
  \item \textbf{SCC}: our SPPFCSPC-CBAM module (Section~\ref{sec:method-sppf})
  in isolation, all other stages unchanged from YOLO11.
  \item \textbf{T}: TOSA + P2 head (Section~\ref{sec:method-tosa}) in
  isolation, SPPF unchanged.
  \item \textbf{COMBO}: our full proposed architecture
  (Section~\ref{sec:method-combined}).
\end{itemize}

\subsection{Training protocol}
All models are trained from the corresponding YOLO11n/YOLO12n
COCO-pretrained checkpoint, for \todo{200} epochs at \todo{1280}px, batch
size 16 (\todo{except COMBO, trained at batch 8 on VisDrone due to a
GPU memory constraint -- confirm which runs this applied to}), with SGD
(lr0$=0.01$, lrf$=0.01$, momentum$=0.937$, weight decay$=5\times10^{-4}$),
a cosine learning-rate schedule, and the standard YOLO detection loss
(box$=7.5$, cls$=0.5$, dfl$=1.5$). All variants are trained with identical
hyperparameters to isolate the effect of the architectural change from any
optimization confound, with the batch-size exception above noted as a
deviation.

\paragraph{Weight initialization.} Pretrained weights are transferred
tensor-by-tensor, matched by exact parameter name and shape between the
COCO checkpoint and the target architecture; any tensor that does not
match keeps PyTorch's default (random) initialization. Because layers are
named by their sequential position in the network, how much of a variant
actually inherits COCO-pretrained weights depends on \emph{where} its
architectural changes sit, not merely on whether \texttt{pretrained=True}
was set. We measured the exact transferred/total tensor count for every
variant (Table~\ref{tab:init}): single-layer, same-position replacements
(GL, SCC) leave every other layer's naming and shape unchanged, so
$>$92\% of tensors still transfer; variants that insert additional layers
early in the backbone (T, COMBO) or replace a block type network-wide
(GLCCP, CCP) shift or rename most subsequent layers, so only
$\approx$8--10\% of tensors transfer -- in practice close to training from
scratch beyond the first stem convolution(s), despite
\texttt{pretrained=True}. We report this because it is a potential
confound for the comparison: part of an observed accuracy gap could stem
from this initialization difference rather than from the architectural
change under study. \todo{Decide whether to control for this (e.g. report
both pretrained and from-scratch runs for the heavily-reindexed variants)
or report it as a limitation.}

\begin{table}[htbp]
\caption{Fraction of model tensors initialized from the COCO-pretrained
checkpoint (vs. left at PyTorch default init), scale n}
\label{tab:init}
\centering
\begin{tabular}{lccc}
\toprule
Variant & Transferred & Total & \% \\
\midrule
YOLO11 (baseline) & 499 & 499  & 100.0 \\
GL                 & 487 & 511  & 95.3 \\
SCC (ours)         & 493 & 532  & 92.7 \\
GLCCP               & 161 & 1749 & 9.2 \\
CCP                 & 173 & 1745 & 9.9 \\
T (ours)           & 56  & 611  & 9.2 \\
COMBO (ours)       & 56  & 665  & 8.4 \\
\bottomrule
\end{tabular}
\end{table}

\subsection{Evaluation metrics}
We report precision (P), recall (R), mAP\textsubscript{50}, and
mAP\textsubscript{50-95}, computed on the validation split with a
confidence threshold of $0.001$ and an NMS IoU threshold of $0.7$, matching
the YOLO validation protocol so that results are directly comparable
across variants. \todo{Add a small-object-only breakdown (e.g. COCO-style
area$<32^2$px bucket) once per-size metrics are available, since aggregate
mAP can dilute a small-object-targeted improvement.}

\section{Planned Large-Scale Results}
\label{sec:results}

The tables in this section report the metrics the forthcoming
experimental phase will populate, in the exact form the final submission
will use; they are deliberately included empty (rather than omitted) so
that the comparison protocol -- which models, which metrics, which
datasets -- is reviewable now, ahead of the results themselves.
\todo{All tables below are placeholders pending completion of full-length
(200-epoch) training runs for every model/dataset pair. Do not submit with
placeholder values.}

\begin{table}[htbp]
\caption{DOTA v1.5 -- validation results}
\label{tab:dota}
\centering
\begin{tabular}{lcccc}
\toprule
Model & P & R & mAP\textsubscript{50} & mAP\textsubscript{50-95} \\
\midrule
YOLO11      & \todo{} & \todo{} & \todo{} & \todo{} \\
YOLO12      & \todo{} & \todo{} & \todo{} & \todo{} \\
GL          & \todo{} & \todo{} & \todo{} & \todo{} \\
GLCCP       & \todo{} & \todo{} & \todo{} & \todo{} \\
CCP         & \todo{} & \todo{} & \todo{} & \todo{} \\
SCC (ours)  & \todo{} & \todo{} & \todo{} & \todo{} \\
T (ours)    & \todo{} & \todo{} & \todo{} & \todo{} \\
COMBO (ours)& \todo{} & \todo{} & \todo{} & \todo{} \\
\bottomrule
\end{tabular}
\end{table}

\begin{table}[htbp]
\caption{VisDrone -- validation results}
\label{tab:visdrone}
\centering
\begin{tabular}{lcccc}
\toprule
Model & P & R & mAP\textsubscript{50} & mAP\textsubscript{50-95} \\
\midrule
YOLO11      & \todo{} & \todo{} & \todo{} & \todo{} \\
YOLO12      & \todo{} & \todo{} & \todo{} & \todo{} \\
GL          & \todo{} & \todo{} & \todo{} & \todo{} \\
GLCCP       & \todo{} & \todo{} & \todo{} & \todo{} \\
CCP         & \todo{} & \todo{} & \todo{} & \todo{} \\
SCC (ours)  & \todo{} & \todo{} & \todo{} & \todo{} \\
T (ours)    & \todo{} & \todo{} & \todo{} & \todo{} \\
COMBO (ours)& \todo{} & \todo{} & \todo{} & \todo{} \\
\bottomrule
\end{tabular}
\end{table}

\begin{table}[htbp]
\caption{Model size and compute cost (scale n)}
\label{tab:cost}
\centering
\begin{tabular}{lcc}
\toprule
Model & Params (M) & GFLOPs \\
\midrule
YOLO11      & 2.59 & \todo{} \\
GLCCP (CCP) & 2.79 & \todo{} \\
SCC (ours)  & 4.27 & \todo{} \\
T (ours)    & 2.49 & \todo{} \\
COMBO (ours)& 4.37 & \todo{} \\
\bottomrule
\end{tabular}
\end{table}

\subsection{Ablation discussion}
\todo{Once Tables~\ref{tab:dota}--\ref{tab:visdrone} are populated, discuss:
(i) whether SCC alone improves over GL/CCP, isolating the contribution of
the CSP shortcut and CBAM attention from the dilation-based context
aggregation used by GL/CCP; (ii) whether T alone (early-stage, P2 head)
outperforms every P5-only variant, as argued in
Section~\ref{sec:intro} on resolution grounds; (iii) whether COMBO's
improvement over the best single-intervention variant is closer to additive
or to redundant, which is the central empirical question this ablation is
designed to answer.}

\section{Conclusion}
\label{sec:conclusion}

We presented a combined architecture for small-object detection in
YOLO11/YOLO12 -- an attention-refined SPPF replacement at the network's
lowest-resolution stage, a high-resolution early-stage attention branch
with a dedicated P2 detection head, and content-aware FPN upsampling --
and, rather than reporting benchmark numbers before the mechanisms behind
them are established, we validated the reasoning behind each component
directly. Grounded in the real annotation statistics of DOTA v1.5 and
VisDrone, every small object in both datasets collapses to under one
feature-map cell by the time it reaches P5, which is why a P5-only
intervention cannot be sufficient on its own and motivates acting at P2
instead. An exact gradient analysis confirms that SPPF's pooling cascade
has identically zero gradient at the 94\% of positions that are not a
local maximum, while the CSP shortcut our SPPFCSPC-CBAM module adds has
nonzero gradient everywhere -- a structural, not merely empirical,
difference. A controlled reconstruction proxy shows content-aware
upsampling recovering 1.17\,dB more PSNR than nearest-neighbor
interpolation after a training budget of minutes, not hours. Together,
these results establish that each intervention addresses a real,
measurable limitation of the baseline through the specific mechanism
claimed for it, ahead of the large-scale comparative training study
(Section~\ref{sec:experiments}) that will determine the magnitude of
their effect on standard detection metrics. Future work includes that
large-scale study itself, extending the analysis to inference-time
techniques orthogonal to architecture (tiled/sliced inference,
small-object-aware NMS, and IoU-loss variants less sensitive to
small-box localization error), and evaluating transfer to additional
aerial benchmarks beyond DOTA and VisDrone.

\section*{Acknowledgments}
\todo{Funding / acknowledgments, if any.}

\bibliographystyle{IEEEtran}
\begin{thebibliography}{00}

\bibitem{redmon2016yolo}
J. Redmon, S. Divvala, R. Girshick, and A. Farhadi, ``You only look once: Unified, real-time object detection,'' in \emph{Proc. IEEE Conf. Comput. Vis. Pattern Recognit. (CVPR)}, 2016.

\bibitem{he2015spp}
K. He, X. Zhang, S. Ren, and J. Sun, ``Spatial pyramid pooling in deep convolutional networks for visual recognition,'' \emph{IEEE Trans. Pattern Anal. Mach. Intell.}, vol. 37, no. 9, pp. 1904--1916, 2015.

\bibitem{wang2023yolov7}
C.-Y. Wang, A. Bochkovskiy, and H.-Y. M. Liao, ``YOLOv7: Trainable bag-of-freebies sets new state-of-the-art for real-time object detectors,'' in \emph{Proc. IEEE/CVF Conf. Comput. Vis. Pattern Recognit. (CVPR)}, 2023.

\bibitem{liu2018rfb}
S. Liu, D. Huang, and Y. Wang, ``Receptive field block net for accurate and fast object detection,'' in \emph{Proc. Eur. Conf. Comput. Vis. (ECCV)}, 2018.

\bibitem{chen2017deeplab}
L.-C. Chen, G. Papandreou, F. Schroff, and H. Adam, ``Rethinking atrous convolution for semantic image segmentation,'' \emph{arXiv preprint arXiv:1706.05587}, 2017.

\bibitem{woo2018cbam}
S. Woo, J. Park, J.-Y. Lee, and I. S. Kweon, ``CBAM: Convolutional block attention module,'' in \emph{Proc. Eur. Conf. Comput. Vis. (ECCV)}, 2018.

\bibitem{yang2021simam}
L. Yang, R.-Y. Zhang, L. Li, and X. Xie, ``SimAM: A simple, parameter-free attention module for convolutional neural networks,'' in \emph{Proc. Int. Conf. Mach. Learn. (ICML)}, 2021.

\bibitem{wang2019carafe}
J. Wang, K. Chen, R. Xu, Z. Liu, C. C. Loy, and D. Lin, ``CARAFE: Content-aware reassembly of features,'' in \emph{Proc. IEEE/CVF Int. Conf. Comput. Vis. (ICCV)}, 2019.

\bibitem{xia2018dota}
G.-S. Xia, X. Bai, J. Ding, Z. Zhu, S. Belongie, J. Luo, M. Datcu, M. Pelillo, and L. Zhang, ``DOTA: A large-scale dataset for object detection in aerial images,'' in \emph{Proc. IEEE Conf. Comput. Vis. Pattern Recognit. (CVPR)}, 2018.

\bibitem{zhu2021visdrone}
P. Zhu, L. Wen, D. Du, X. Bian, H. Fan, Q. Hu, and H. Ling, ``Detection and tracking meet drones challenge,'' \emph{IEEE Trans. Pattern Anal. Mach. Intell.}, vol. 44, no. 11, pp. 7380--7399, 2021.

\end{thebibliography}

\end{document}
